\documentclass{article}
\usepackage[utf8]{inputenc}
\usepackage{geometry}
\geometry{a4paper, margin=1in}

\title{The Story We Pass On: Loops of Teaching and Practice}
\author{Albert Jan van Hoek}
\date{\today}

\begin{document}

\maketitle

What story do I want to tell my children about life?

A story that helps them understand how life works — so they can survive it, improve it, and pass it on.

Teaching that becomes preaching that becomes teaching again.

You can see that as a chain.

Or — more accurately — as a loop.

\section*{The Loop of Collaboration}

This became central in my \textbf{Theory of Long-term Collaboration}.
Collaboration between learning systems becomes stable when the behaviours that improve collaboration reinforce themselves over time.

We teach what works.
We practise what we teach.
We normalise what we practise.
And that shapes the next round of teaching.

When preaching and practice align, systems sustain themselves.

\section*{The Power of Hypocrisy}

However, how to dissolve a functional loop?

When practice and preaching diverge, the loops fracture.

Such fracture has a name: \textbf{hypocrisy}.

And hypocrisy is structurally powerful.

It can quietly dismantle a healthy loop:
\begin{itemize}
    \item ``Make the nation great again'' — while increasing inequality.
    \item ``Defend freedom'' — while narrowing who benefits from it.
    \item ``Support working people'' — while cutting taxes for the rich.
    \item ``Protect the future'' — while accelerating short-term extraction.
\end{itemize}

The slogan points one way.
The policy reinforces another.

The loop follows behaviour — not the banner.

\section*{Inversion as Repair}

But here is where it becomes interesting.
If hypocrisy can erode, it can also expose.

Call it hypocrisy squared.
\begin{itemize}
    \item ``Make the nation worse again'' — while decreasing inequality.
    \item ``Choose autocracy'' — while expanding democratic participation.
    \item ``Stand with billionaires'' — while implementing progressive taxation.
    \item ``We aim to be the last generation'' — while fully committing to climate agreements.
\end{itemize}

The words sound absurd.
The actions repair the loop.

That is the inversion.
If someone uses the language of virtue to reinforce harm,
you can use the language of vice to reinforce repair.

From inside the old loop, that looks insane.
From outside, it reveals the structure.

\section*{The Lesson}

Because systems don’t run on slogans.
They run on reinforcement.

So perhaps that is the story worth passing on:
\begin{itemize}
    \item Don’t listen to the banner.
    \item Watch what is rewarded.
    \item Align what you teach with what you practise.
    \item And when rhetoric hides reinforcement in the wrong direction — invert it.
\end{itemize}

\end{document}
